\begin{tikzpicture}[
    font=\sffamily\scriptsize,
    >={Latex[length=2.4mm]},
    flow/.style={->,very thick,draw=black!70},
    labelbox/.style={fill=white,draw=black!20,rounded corners=1.5pt,
        inner xsep=3pt,inner ysep=2pt,align=center},
    store/.style={draw=black!65,rounded corners=2pt,align=center,
        minimum height=12mm,fill=blue!6},
    process/.style={draw=black!65,rounded corners=2pt,align=center,
        minimum height=12mm,fill=orange!10}
]
% Panel A: strip-wise WSI access.
\node[font=\sffamily\small\bfseries,anchor=west] at (0,4.7)
    {(a) Lectura espacial de una WSI};
\draw[rounded corners=2pt,fill=black!3,draw=black!55] (0,0.4) rectangle (5.1,4.35);
\path[fill=pink!30,draw=purple!45!black,line width=0.6pt]
    (0.55,1.05) .. controls (0.15,2.0) and (0.65,3.65) .. (1.8,3.95)
    .. controls (2.75,4.2) and (3.15,3.25) .. (4.35,3.45)
    .. controls (5.0,3.0) and (4.55,2.35) .. (4.8,1.65)
    .. controls (4.1,0.75) and (3.1,1.35) .. (2.35,0.85)
    .. controls (1.65,0.4) and (1.15,0.95) .. cycle;
\foreach \x in {0.85,1.55,2.25,2.95,3.65,4.35}
    \draw[black!18] (\x,0.4) -- (\x,4.35);
\foreach \y in {1.1,1.8,2.5,3.2,3.9}
    \draw[black!18] (0,\y) -- (5.1,\y);
\fill[blue!35,opacity=0.32] (0.03,1.82) rectangle (5.07,2.48);
\draw[blue!65!black,line width=1.1pt] (0.03,1.82) rectangle (5.07,2.48);
\foreach \x in {0.85,2.25,3.65}
    \draw[orange!90!black,line width=1.5pt] (\x,1.8) rectangle ++(0.7,0.7);
\node[labelbox,anchor=south west] at (0.12,0.52)
    {cuadrícula sin solapamiento\\celda $256\times256$};

\draw[fill=blue!10,draw=blue!60!black,rounded corners=1pt]
    (7.1,1.65) rectangle (12.9,2.65);
\foreach \x in {7.25,8.15,9.05,9.95,10.85,11.75}
    \draw[black!25] (\x,1.65) rectangle ++(0.9,1.0);
\foreach \x in {7.25,9.05,10.85}
    \draw[orange!90!black,line width=1.5pt] (\x,1.65) rectangle ++(0.9,1.0);
\node[labelbox,above=3mm] at (10.0,2.65)
    {recorte de las posiciones seleccionadas\\sin decodificar la WSI completa en memoria};
% El rótulo se dibuja después de la franja para conservar todo el texto visible.
\draw[flow] (5.15,2.15) -- node[labelbox,above=1mm]
    {una franja horizontal\\de 256 píxeles de alto} (7.05,2.15);

% Panel B: the two storage routes.
\node[font=\sffamily\small\bfseries,anchor=west] at (0,-0.25)
    {(b) Rutas de almacenamiento};
\node[process,minimum width=31mm] (trainpatch) at (1.7,-1.4)
    {Parches RGB de\\preentrenamiento};
\node[store,minimum width=40mm,right=14mm of trainpatch] (bin)
    {archivo prebarajado \texttt{.bin}\\cabecera de 64 bytes\\registros CHW \texttt{uint8}};
\node[process,minimum width=33mm,right=14mm of bin] (mmap)
    {\texttt{mmap} de solo lectura\\índice $\mapsto$ desplazamiento\\vista sin copia};
\draw[flow] (trainpatch) -- node[labelbox,above=6mm]{orden fijado} (bin);
\draw[flow] (bin) -- node[labelbox,above=6mm]{lectura secuencial} (mmap);

\node[process,minimum width=31mm] (evalpatch) at (1.7,-3.25)
    {Franjas RGB de la\\cohorte de evaluación};
\node[store,minimum width=40mm,right=14mm of evalpatch] (enc)
    {dos codificadores congelados\\mismas coordenadas};
\node[process,minimum width=33mm,right=14mm of enc] (latent)
    {almacenes FP32 pareados\\$\mu[16,32,32]$\\sin copia RGB intermedia};
\draw[flow] (evalpatch) -- node[labelbox,above=7mm]{cada parche una vez} (enc);
\draw[flow] (enc) -- node[labelbox,above=7mm]{escritura directa} (latent);
\end{tikzpicture}
